\subsubsection{Predictive power by bonus rule}

In this subsection, we briefly analyze whether the predictive power of the strategic \aisub{} is sensitive to the different types of games from Table~\ref{tab:game_parameters} (and therefore different possible distributions for $\pi$).
In particular, we focus on the bonus rules, although the results are similar for all dimensions of the parameter space.
As a reminder, different bonus rules change the fundamental structure of the game.
For some, many strategy profiles led to the bonus (e.g., ``Unequal''), and others only one (``Coord. Low'').
Broadly speaking, the rules can be categorized into two types: mutually achievable and competitive.
In mutually achievable games, both players receive the bonus when the bonus condition is met.
In competitive games, only one player can receive the bonus.
This distinction roughly corresponds to the classic game-theoretic categories of coordination versus zero-sum games.

Figure~\ref{fig:bonus_rule_prop_plot} shows the proportion of games for which the strategic \aisub{} is the best predictor of initial play by bonus rule ($\sum_{s \in B} \mathbf 1\{\hat \Lambda_s > 0\} / |B|$ where $B \subset S$ is the set of games with a given bonus rule).
The figure is of a similar structure as the bottom panel of Figure~\ref{fig:ll_ratios}, with each panel showing a different benchmark comparison with $\varepsilon = 0.2$.
The y-axis shows the bonus rule, and the x-axis is the proportion of games for which the log-likelihood ratio is greater than zero.
Green indicates that strategic \aisub{} significantly outperforms the reference model in more than 50\% of games with that bonus rule, and grey indicates no significant difference.
Bonus rules above the horizontal solid black line correspond to games where both players can achieve the bonus simultaneously, while those below are competitive, permitting only one player to do so.

\begin{figure}[h!]
\begin{center}
\begin{minipage}{\textwidth}
\caption{Relative predictive power of strategic \aisub{s} ($\varepsilon = 0.2$) by bonus rule} 
\label{fig:bonus_rule_prop_plot}
\vspace*{-0.15in}
\includegraphics[width=\textwidth]{plots/bonus_rule_prop_plot_eps_02.pdf} 
\end{minipage}
\end{center}
\begin{footnotesize}
\begin{singlespace}
\vspace*{-0.05in}
\emph{Notes:} This figure shows the proportion of games for which the strategic \aisub{s} is the best predictor of initial play for each bonus rule (on the y-axis).
The vertical dashed line corresponds to a 50-50 split, where there is no difference between the strategic \aisub{} and the reference model in that panel.
Bonus rules above the horizontal solid black line are mutually achievable (both players can receive the bonus), and those below are competitive (only one player can receive the bonus).
Green indicates that strategic \aisub{} significantly outperforms the reference model in more than 50\% of games with that bonus rule, and grey indicates no significant difference. 
Error bars show 95\% Clopper-Pearson confidence intervals.
The bonus rules are as follows:
\textbf{Gap Low}---selecting exactly \{\textit{gap}\} less than the opponent;
\textbf{Gap High}---selecting exactly \{\textit{gap}\}  more than the opponent;
\textbf{Gap Abs.}---when absolute difference equals \{\textit{gap}\} ;
\textbf{More Than}---when difference exceeds \{\textit{gap}\} ;
\textbf{Equal}---for matching opponent's number;
\textbf{Unequal}---for selecting different number than opponent;
\textbf{Sum Even}---when sum of both numbers is even;
\textbf{Sum Odd}---when sum of both numbers is odd;
\textbf{Sum Upper}---when sum equals the upper bound;
\textbf{Less Upper}---when the sum is less than the upper bound;
\textbf{Coord. Low}---when both players select the lower bound.
\end{singlespace}
\end{footnotesize}
\end{figure}

The strategic \aisub{} sample weakly dominates all other benchmarks; it never significantly underperforms for any bonus rule.
It outperforms the baseline and the random pure-strategy benchmark for every rule, the cognitive hierarchy model and the Harsanyi-Selten Nash equilibria for all but two rules, and the uniform benchmark for 7 of the 11 rules.
There is no apparent distinction between mutually achievable and competitive bonus rules.

In Appendix~\ref{app:robustness_checks}, we provide analogous plots for alternative values of $\varepsilon \in \{0.05, 0.1, 0.3\}$ and for the various points rules in Table~\ref{tab:game_parameters}.
We also report regressions of other game parameters (lower bound, upper bound, bonus size, gap) on the log-likelihood ratio of the strategic \aisub{} relative to the benchmarks ($\hat \Lambda_s$).
The results are similar to the ones presented here.
The strategic \aisub{s} are almost universally superior across different slices of the parameter space.
They are therefore likely to generalize to alternative distributions for $\pi$ across the population of games.
